% =============================================================================
% preamble.tex — giao diện chung của báo cáo: gói, màu, hộp, macro.
% Kế thừa phong cách của mẫu báo cáo xv6 Lab 01 (XeLaTeX, DejaVu, tcolorbox).
% Khi viết nội dung KHÔNG cần sửa file này. Cần thêm gói mới thì báo người ghép
% báo cáo (TV4) để cả nhóm dùng chung một bản.
% =============================================================================

% --- Trang và lề -------------------------------------------------------------
\usepackage[a4paper,margin=2.2cm,top=2.5cm,bottom=2.5cm]{geometry}

% --- Phông chữ Unicode tiếng Việt --------------------------------------------
% Dùng họ DejaVu (có sẵn trên Overleaf). Máy nào thiếu thì tự lùi về Latin Modern
% để vẫn biên dịch được; khi đó dấu tiếng Việt có thể kém đẹp hơn.
% Phông mã dùng DejaVu Sans, KHÔNG dùng DejaVu Sans Mono: bản Mono thiếu một số chữ
% tiếng Việt có hai dấu (ấ, ế, ố, ầ, ở, ỉ), nên mã giả và chú thích tiếng Việt bị mất
% ký tự. Mẫu báo cáo xv6 cũng chọn như vậy.
\usepackage{fontspec}
\IfFontExistsTF{DejaVu Serif}%
  {\setmainfont{DejaVu Serif}}%
  {\setmainfont{Latin Modern Roman}}
\IfFontExistsTF{DejaVu Sans}%
  {\setsansfont{DejaVu Sans}\setmonofont{DejaVu Sans}}%
  {\setsansfont{Latin Modern Sans}\setmonofont{Latin Modern Mono}}

\usepackage{polyglossia}
\setdefaultlanguage{vietnamese}

% --- Toán, đồ họa, bảng, danh sách -------------------------------------------
\usepackage{amsmath,amssymb,amsthm}
\usepackage{etoolbox}
\usepackage{xparse}
\usepackage{float}        % cho phép [H]: bảng và hình nằm đúng chỗ khai báo
\usepackage{graphicx}
\usepackage{booktabs}
\usepackage{array}
\usepackage{tabularx}
\usepackage{longtable}
\usepackage{calc}
\usepackage{fvextra}
\usepackage{multirow}
\usepackage{enumitem}
\usepackage[table]{xcolor}
\usepackage{caption}
\usepackage{hyperref}
\usepackage{xurl}
\usepackage{pdflscape}
\usepackage{fancyhdr}
\usepackage{titlesec}
\usepackage{tcolorbox}
\usepackage{fancyvrb}
\usepackage{lastpage}
\usepackage{microtype}
\usepackage{tikz}
\tcbuselibrary{skins,breakable}
\usetikzlibrary{arrows.meta,positioning,calc}

% --- Bảng màu (giữ nguyên bảng màu của mẫu xv6) ------------------------------
\definecolor{primary}{RGB}{20, 60, 105}        % xanh học thuật đậm
\definecolor{secondary}{RGB}{41, 128, 185}     % xanh biển hiện đại
\definecolor{darkneutral}{RGB}{44, 62, 80}     % xám đen cho chữ
\definecolor{lightbg}{RGB}{247, 249, 252}      % nền nhạt cho hộp và mã
\definecolor{tableheader}{RGB}{232, 241, 250}  % nền tiêu đề bảng
\definecolor{codeborder}{RGB}{210, 222, 235}   % viền hộp
\definecolor{accent}{RGB}{231, 76, 60}         % màu nhấn: chỗ chưa điền, cảnh báo
% Màu các loại ô trong hình minh họa bàn chơi
\definecolor{tilefloor}{RGB}{200, 214, 232}
\definecolor{tilefragile}{RGB}{245, 166, 35}
\definecolor{tilebridge}{RGB}{222, 190, 120}

% --- Liên kết PDF ------------------------------------------------------------
% \reportpdftitle và \reportpdfauthor được khai báo ở đầu Report.tex.
\hypersetup{
  colorlinks=true,
  linkcolor=primary,
  urlcolor=secondary,
  citecolor=primary,
  unicode=true,
  bookmarksnumbered=true,
  pdftitle={\reportpdftitle},
  pdfauthor={\reportpdfauthor}
}

% --- Khoảng cách dòng, đoạn, danh sách ---------------------------------------
\linespread{1.18}
\setlength{\parindent}{0pt}
\setlength{\parskip}{0.55em}
\setlength{\emergencystretch}{2em}   % bớt dòng tràn lề khi có nhiều tên hàm không ngắt được
\setlist[itemize]{leftmargin=1.6em,itemsep=0.2em,topsep=0.2em}
\setlist[enumerate]{leftmargin=1.8em,itemsep=0.2em,topsep=0.2em}
\setcounter{tocdepth}{2}      % mục lục tới mục con 3.1
\setcounter{secnumdepth}{-2}  % số mục đã viết sẵn trong tiêu đề (2.1, 3.4.2, ...)

% Hình, bảng, công thức đánh số theo chương: Hình 5.1, Bảng 5.2 (docs/06 §2)
\numberwithin{equation}{section}
\numberwithin{figure}{section}
\numberwithin{table}{section}

\captionsetup{font=small,labelformat=empty,labelsep=none,
  justification=justified,singlelinecheck=true}
\captionsetup[table]{position=top,skip=5pt}
\captionsetup[figure]{position=bottom,skip=6pt}

% --- Tiêu đề các mục (titlesec) ----------------------------------------------
\titleformat{\section}
  {\color{primary}\Large\bfseries}
  {\thesection.}{0.6em}{}[\vspace{0.2em}{\color{primary}\titlerule[1.2pt]}]
\titleformat{\subsection}
  {\color{secondary}\large\bfseries}
  {\thesubsection}{0.6em}{}
\titleformat{\subsubsection}
  {\color{darkneutral}\normalsize\bfseries}
  {\thesubsubsection}{0.5em}{}

% Mỗi chương (\section) bắt đầu ở trang mới: tiêu đề không bị bỏ lẻ cuối trang và mỗi
% người viết một chương không làm xô lệch trang của chương khác.
\pretocmd{\section}{\clearpage\phantomsection}{}{}

% --- Công tắc bản nháp ---------------------------------------------------------
% Bản nháp: hiện hộp hướng dẫn viết và chữ BẢN NHÁP ở đầu trang.
% Bản nộp: đặt \bannhapfalse trong Report.tex rồi biên dịch lại.
\newif\ifbannhap
\bannhapfalse

% --- Đầu trang và chân trang -------------------------------------------------
% \headerleft và \headerright được khai báo ở đầu Report.tex.
\pagestyle{fancy}
\fancyhf{}
\lhead{\small\textcolor{darkneutral}{\textbf{\headerleft}}}
\chead{\ifbannhap\small\textcolor{accent}{\textbf{BẢN NHÁP}}\fi}
\rhead{\small\textcolor{gray}{\headerright}}
\rfoot{\small\textcolor{darkneutral}{\textbf{Trang \thepage\,/\,\pageref{LastPage}}}}
\renewcommand{\headrulewidth}{0.6pt}
\renewcommand{\headrule}{\hbox to\headwidth{\color{secondary!70}\leaders\hrule height \headrulewidth\hfill}}
\renewcommand{\footrulewidth}{0.4pt}
\renewcommand{\footrule}{\hbox to\headwidth{\color{gray!30}\leaders\hrule height \footrulewidth\hfill}}
\setlength{\headheight}{27pt}

\AtBeginDocument{%
  \ifbannhap
    \PackageWarningNoLine{bloxreport}{Dang o che do BAN NHAP (\string\bannhaptrue). Dat \string\bannhapfalse\space trong Report.tex khi xuat ban de nop}%
  \fi}

% --- Macro đánh dấu chỗ chưa viết -------------------------------------------
% \placeholder{...}  chỗ cần điền nội dung (chữ đỏ nghiêng trong ngoặc vuông)
% \choTV{TV3}{...}   phần người khác phải viết, không tự viết thay
\newcommand{\placeholder}[1]{\textcolor{accent!85!black}{\textit{[#1]}}}
\newcommand{\choTV}[2]{\textcolor{accent!85!black}{\textbf{[#1 viết:}\;\textit{#2}\textbf{]}}}

% Tên hàm, biến, tệp trong mã nguồn: \code{open_bridges} (không cần gõ \_)
\newcommand{\code}[1]{\texttt{\detokenize{#1}}}
% Từ khóa trong mã giả: \kw{nếu}
\newcommand{\kw}[1]{\textbf{#1}}

% --- Tham chiếu có liên kết: \bang{tab:x} -> "Bảng 2.1" ------------------------
\newcommand{\bang}[1]{\hyperref[#1]{Bảng~\ref*{#1}}}
\newcommand{\hinh}[1]{\hyperref[#1]{Hình~\ref*{#1}}}
\newcommand{\chuong}[1]{\hyperref[#1]{chương~\ref*{#1}}}
\newcommand{\muc}[1]{\hyperref[#1]{mục~\ref*{#1}}}

% --- Hộp hiển thị mã, ghi chú -------------------------------------------------
\newtcolorbox{terminalbox}[1][]{
  enhanced, colback=lightbg, colframe=codeborder, coltitle=primary,
  boxrule=0.7pt, arc=3.5pt,
  left=10pt,right=10pt,top=8pt,bottom=8pt,
  breakable, #1
}

\newtcolorbox{notebox}[1][]{
  enhanced, colback=primary!4, colframe=primary,
  boxrule=0pt, leftrule=3.5pt, arc=2pt,
  left=10pt,right=10pt,top=6pt,bottom=6pt,
  breakable, #1
}

% Hộp mã giả theo phong cách thuật toán trong bài báo khoa học
\newenvironment{algorithmblock}[1]{%
  \begin{tcolorbox}[
    enhanced, breakable, colback=white, colframe=primary!55,
    colbacktitle=primary, coltitle=white,
    boxrule=0.7pt, arc=2pt,
    left=10pt,right=10pt,top=7pt,bottom=8pt,
    title={#1}, fonttitle=\small
  ]
}{%
  \end{tcolorbox}
}

% Trong algorithmcode: dùng \kw{...} cho từ khóa. Dấu ngoặc nhọn và dấu gạch chéo
% ngược là ký tự đặc biệt; tránh chúng trong mã giả (dùng ngoặc tròn hoặc chữ).
\DefineVerbatimEnvironment{algorithmcode}{Verbatim}{
  commandchars=\\\{\},
  fontsize=\small,
  baselinestretch=1.05,
  numbers=left,
  numbersep=8pt,
  xleftmargin=2.0em,
  frame=none
}

% --- Hộp hướng dẫn viết (chỉ hiện khi \bannhaptrue) ------------------------------
% Dùng: \guide{ nội dung hướng dẫn }. Không đặt môi trường verbatim bên trong.
\newtcolorbox{guidebox}{
  enhanced, breakable,
  colback=yellow!10, colframe=orange!70!black, colbacktitle=orange!25, coltitle=black,
  boxrule=0.6pt, arc=2pt,
  left=8pt,right=8pt,top=5pt,bottom=5pt,
  fonttitle=\small\bfseries,
  title={Hướng dẫn viết \textnormal{(chỉ hiện ở bản nháp, xóa khi nộp)}}
}
\newcommand{\guide}[1]{\ifbannhap\begin{guidebox}\small\raggedright #1\end{guidebox}\fi}

% --- Hình có khung giữ chỗ khi chưa có tệp ảnh ---------------------------------
% \hinhanh{tên-tệp.png}{chiều rộng}{chú thích}{nhãn}
% Tệp đặt trong figures/. Chưa có tệp thì hiện khung giữ chỗ, báo cáo vẫn biên dịch được.
\newcommand{\figplaceholder}[1]{%
  \begin{tcolorbox}[enhanced, colback=lightbg, colframe=accent!60, boxrule=0.8pt,
    frame style={dashed}, arc=2pt, width=0.92\linewidth,
    halign=center, valign=center, height=4.2cm]
    \textcolor{accent!85!black}{\textit{[Chèn hình: \texttt{figures/\detokenize{#1}}]}}
  \end{tcolorbox}}
\newcommand{\hinhanh}[4]{%
  \begin{figure}[H]
    \centering
    \IfFileExists{figures/#1}%
      {\includegraphics[width=#2]{figures/#1}}%
      {\figplaceholder{#1}}
    \caption{#3}\label{#4}
  \end{figure}}

% --- Định lý, định nghĩa, ví dụ (đánh số theo chương) -------------------------
\theoremstyle{definition}
\newtheorem{dinhnghia}{Định nghĩa}[section]
\newtheorem{vidu}{Ví dụ}[section]
\theoremstyle{plain}
\newtheorem{dinhly}{Định lý}[section]
\newtheorem{bode}{Bổ đề}[section]
\newtheorem{menhde}{Mệnh đề}[section]
\theoremstyle{remark}
\newtheorem{nhanxet}{Nhận xét}[section]
\renewcommand{\proofname}{Chứng minh}
\gappto\captionsvietnamese{\renewcommand{\proofname}{Chứng minh}\renewcommand{\refname}{Tài liệu tham khảo}}

% --- Kiểu TikZ dùng chung --------------------------------------------------------
\tikzset{
  flowbox/.style={draw=primary!65,fill=lightbg,rounded corners=2pt,
    line width=0.7pt,align=center,font=\small,
    text width=3.7cm,minimum height=1.2cm,inner sep=5pt},
  flowarrow/.style={-{Stealth[length=2.2mm]},line width=0.8pt,draw=primary}
}

% --- Các loại ô của Bloxorz (docs/02 §2) ----------------------------------------
% Mỗi macro \tUnit... vẽ một ô đơn vị (0,0)-(1,1) trong một tikzpicture đang mở.
% Ký hiệu trong file màn: . # S G F o x @ b B
\newcommand{\tUnitVoid}{}
\newcommand{\tUnitFloor}{\path[fill=tilefloor,draw=darkneutral!45,line width=0.3pt] (0,0) rectangle (1,1);}
\newcommand{\tUnitStart}{\tUnitFloor\node[font=\bfseries\tiny,text=primary] at (.5,.5) {S};}
\newcommand{\tUnitGoal}{\tUnitFloor\fill[black!85] (.5,.5) circle (.33);\draw[white,line width=0.3pt] (.5,.5) circle (.2);}
\newcommand{\tUnitFragile}{\path[fill=tilefragile,draw=darkneutral!45,line width=0.3pt] (0,0) rectangle (1,1);\draw[orange!30!black,line width=0.35pt] (.22,1)--(.45,.55)--(.3,.4)--(.62,0);}
\newcommand{\tUnitSoft}{\tUnitFloor\draw[darkneutral,line width=0.6pt] (.5,.5) circle (.27);}
\newcommand{\tUnitHeavy}{\tUnitFloor\draw[darkneutral,line width=0.7pt] (.22,.22)--(.78,.78) (.22,.78)--(.78,.22);}
\newcommand{\tUnitSplit}{\tUnitFloor\draw[darkneutral,line width=0.6pt] (.36,.2) to[bend right=45] (.36,.8) (.64,.2) to[bend left=45] (.64,.8);}
\newcommand{\tUnitBridgeOpen}{\path[fill=tilebridge,draw=darkneutral!45,line width=0.3pt] (0,0) rectangle (1,1);\draw[brown!60!black,line width=0.3pt] (0,.33)--(1,.33) (0,.66)--(1,.66);}
\newcommand{\tUnitBridgeClosed}{\path[fill=white,draw=brown!60!black,dashed,line width=0.4pt] (0,0) rectangle (1,1);\draw[brown!40,line width=0.3pt] (0,.33)--(1,.33) (0,.66)--(1,.66);}

% Ô minh họa nhỏ dùng trong dòng chữ và trong bảng: \tFloor, \tGoal, ...
\newcommand{\tilesw}[1]{\begin{tikzpicture}[baseline={(0,0.2)},x=0.42cm,y=0.42cm]#1\end{tikzpicture}}
\newcommand{\tVoid}{\tilesw{\path[draw=gray!55,dotted,line width=0.4pt] (0,0) rectangle (1,1);}}
\newcommand{\tFloor}{\tilesw{\tUnitFloor}}
\newcommand{\tStart}{\tilesw{\tUnitStart}}
\newcommand{\tGoal}{\tilesw{\tUnitGoal}}
\newcommand{\tFragile}{\tilesw{\tUnitFragile}}
\newcommand{\tSoft}{\tilesw{\tUnitSoft}}
\newcommand{\tHeavy}{\tilesw{\tUnitHeavy}}
\newcommand{\tSplit}{\tilesw{\tUnitSplit}}
\newcommand{\tBridgeOpen}{\tilesw{\tUnitBridgeOpen}}
\newcommand{\tBridgeClosed}{\tilesw{\tUnitBridgeClosed}}

% --- Vẽ bàn chơi từ dữ liệu màn ---------------------------------------------------
% \bloxboard[tùy chọn tikz]{hàng1/hàng2/...}
% \bloxboardplus[tùy chọn tikz]{hàng1/hàng2/...}{mã tikz vẽ thêm}
% Mỗi hàng là một dòng `grid` trong file màn (Source/levels/original/NN.json), các
% hàng cách nhau bằng dấu `/`, không có dấu cách. Ký tự lạ được vẽ như sàn.
% Trong mã vẽ thêm, ô (hàng r, cột c) có tâm tại tọa độ (c+0.5, -(r+0.5)).
% Không dùng bên trong đối số của macro khác (ví dụ \guide{...}).
\ExplSyntaxOn
\int_new:N \l__blox_r_int
\int_new:N \l__blox_c_int
\tl_new:N \l__blox_x_tl
\tl_new:N \l__blox_y_tl
\seq_new:N \l__blox_rows_seq

\NewDocumentCommand{\bloxdrawtile}{m}
 {
  \str_case:nnF {#1}
   {
    {.}{\tUnitVoid}
    {S}{\tUnitStart}
    {G}{\tUnitGoal}
    {F}{\tUnitFragile}
    {o}{\tUnitSoft}
    {x}{\tUnitHeavy}
    {@}{\tUnitSplit}
    {b}{\tUnitBridgeClosed}
    {B}{\tUnitBridgeOpen}
   }
   {\tUnitFloor}
 }

\cs_new_protected:Npn \__blox_picture:nnn #1#2#3
 {
  \begin{tikzpicture}[#1,x=0.4cm,y=0.4cm]
   \seq_set_split:Nnn \l__blox_rows_seq { / } { #2 }
   \int_zero:N \l__blox_r_int
   \seq_map_inline:Nn \l__blox_rows_seq
    {
     \int_zero:N \l__blox_c_int
     \str_map_inline:nn { ##1 }
      {
       \tl_set:Nx \l__blox_x_tl { \int_use:N \l__blox_c_int }
       \tl_set:Nx \l__blox_y_tl { \int_eval:n { - \l__blox_r_int - 1 } }
       \begin{scope}[shift={(\l__blox_x_tl,\l__blox_y_tl)}]
         \bloxdrawtile { ####1 }
       \end{scope}
       \int_incr:N \l__blox_c_int
      }
     \int_incr:N \l__blox_r_int
    }
   #3
  \end{tikzpicture}
 }

\NewDocumentCommand{\bloxboard}{ O{} v }
 { \__blox_picture:nnn {#1} {#2} {} }
\NewDocumentCommand{\bloxboardplus}{ O{} v m }
 { \__blox_picture:nnn {#1} {#2} {#3} }
\ExplSyntaxOff

% =============================================================================
% Phần thêm cho bản ghép từ Markdown (pandoc)
% =============================================================================
\makeatletter
\newsavebox\pandoc@box
\newcommand*\pandocbounded[1]{%
  \sbox\pandoc@box{#1}%
  \Gscale@div\@tempa{\textheight}{\dimexpr\ht\pandoc@box+\dp\pandoc@box\relax}%
  \Gscale@div\@tempb{\linewidth}{\wd\pandoc@box}%
  \ifdim\@tempb\p@<\@tempa\p@\let\@tempa\@tempb\fi
  \ifdim\@tempa\p@<\p@\scalebox{\@tempa}{\usebox\pandoc@box}%
  \else\usebox{\pandoc@box}%
  \fi}
\makeatother
\newcounter{none} % bảng không có chú thích (pandoc dùng \LTcaptype{none})
\providecommand{\tightlist}{\setlength{\itemsep}{0pt}\setlength{\parskip}{0pt}}

% Hình nằm đúng chỗ khai báo, ngay dưới đoạn nhắc tới nó
\floatplacement{figure}{H}

% Bảng: chữ nhỏ, chú thích rộng bằng dòng, khoảng cách quanh bảng
\setlength{\LTcapwidth}{\linewidth}
\setlength{\LTpre}{8pt}
\setlength{\LTpost}{8pt}
\newcommand{\tablefont}{\small}
\AtBeginEnvironment{longtable}{\tablefont\renewcommand{\arraystretch}{1.12}}
\captionsetup[table]{position=top,skip=4pt,font=small}

% Khối mã và mã giả: nền nhạt, tự xuống dòng khi dòng dài
\RecustomVerbatimEnvironment{verbatim}{Verbatim}{%
  fontsize=\footnotesize, breaklines=true, breakanywhere=true,
  frame=single, framesep=7pt, rulecolor=\color{codeborder},
  formatcom=\color{darkneutral}, baselinestretch=1.05}

% Trích dẫn (> ...) trong Markdown hiện thành hộp ghi chú
\renewenvironment{quote}{\begin{notebox}}{\end{notebox}}


% Ký tự kết thúc chứng minh: DejaVu Serif không có glyph này
\usepackage{newunicodechar}
\newunicodechar{∎}{\ensuremath{\blacksquare}}
